% Transcription — fonds Alexandre Grothendieck, shelfmark 63, batch 1 (pages 1-20).
% Established from the facsimile archives/batches/63/batch-01.pdf (a copy of
% 63.pdf fetched through the project relay, no cover sheet: PDF page k =
% archivists' page k), per the skill transcribe-grothendieck. The folder is a
% hundred pages in five batches; this is the first. Folder 62 (transcripts/62/)
% is the notational reference for the group « Courbes elliptiques » (62-64).
%
% Pass: Opus 5.5 (claude-opus-5-5), 2026-09-26 — first pass, unchecked against the pages by a human.
%
% Pages skipped: 1 and 14, blank cream sheets (nothing written on them).
%
% Pages given one \note each, not re-set (someone else's text):
% - Pages 2-13: a copy (brown, on aged paper) of a HANDWRITTEN MANUSCRIPT BY
%   JOHN TATE, in English: a contents page (p. 2) and ten leaves numbered 1
%   to 10 at the top right (pp. 3-12), plus an eleventh leaf numbered 11
%   (p. 13). Evidence that the hand is not Grothendieck's: a rounded,
%   upright, careful hand with nothing of his fast French hand (compared
%   with folder 62, e.g. tiles of p. 41); English throughout; the text
%   says « See Serre's "Courbes Elliptiques - Formulaire" for the details »
%   (p. 4) and « (Serre tells me) » (p. 8), which the author of the text
%   says, not Grothendieck. Evidence of identity: the contents page lists
%   exactly the eight sections (Generalized Weierstrass form; Change of
%   coordinates; "Minimal" Weierstrass equations over a valuation ring; The
%   canonical filtration on the group of v-adic points; Application: L_v(1)
%   and the integral of |omega_v| when k is finite; The Néron minimum model;
%   Algorithm ... first five cases; ... last five cases) of J. Tate,
%   « Algorithm for determining the type of a singular fiber in an elliptic
%   pencil », Modular Functions of One Variable IV, LNM 476 (1975), 33-52,
%   whose first footnote says it is « a slightly edited version of a letter
%   to Cassels » (checked 2026-09-26 against the scan at
%   wstein.org/Tables/antwerp/tate/tate.pdf; the printed version adds a §0
%   Summary). These leaves are therefore a copy of Tate's letter, a text by
%   another author since published. By the folder rule for other people's
%   texts they get one French \note per page, identifying the content, and
%   are not re-set. No annotation in a second hand is visible on them; the
%   boxes, marginal brackets (« subprocedure branch (7) », « branch 3 »),
%   « (sauf erreur) » (p. 7) and the blacked-out patches (pp. 4, 5, 9, 10)
%   are in the same ink and hand as the text and are Tate's.
% - Pages 15-20: the first six leaves (typed pages 1-6, the first
%   unnumbered, then « - 2 - » ... « - 6 - ») of the typescript « COURBES
%   ELLIPTIQUES : FORMULAIRE d'après J. Tate mis au goût du jour par P.
%   Deligne », §1 « Formules générales » and the start of §2 « Cas où 2 et 3
%   sont inversibles » up to Proposition 2.5. This is Deligne's text,
%   published as P. Deligne, « Courbes elliptiques : formulaire d'après J.
%   Tate », LNM 476 (1975), 53-73 — the title line on p. 15 is the printed
%   one. Not his own text, so one \note per page with the pencil marks in
%   full. The marks are in pencil, in a fast French hand; the only one with
%   words (p. 20, a three-case summary « Sur un corps ... ») looks like his
%   hand but the attribution is not certain, and the notes say so.
%
% His pagination: none of his own on these leaves. Tate's leaves carry
% Tate's numbers 1-11 (recorded per page); the typescript carries typed
% page numbers 1-6. Numbered runs are Tate's (§§1-8, formulae (1.1)-(8.2),
% steps (1)-(10)) and Deligne's ((1.1)-(1.8), (2.1)-(2.4), Prop. 2.5).
% No date on any leaf.
%
% What the batch is about: reference material for folder 63's
% « formulaire » — Tate's generalised Weierstrass equation
% y^2 + a1 xy + a3 y = x^3 + a2 x^2 + a4 x + a6 with b2..b8, c4, c6, Delta,
% j, changes of coordinates, minimal models, the filtration E(K) > E0(K) >
% E1(K) > ..., the local factor L_v(1), Kodaira-Néron types and Tate's
% algorithm; then Deligne's scheme-theoretic reworking over a base S
% (omega = e^* Omega^1, p_* O(ne), the same formulae). Nothing in this batch
% is Grothendieck's mathematics except, possibly, the pencil marks on the
% typescript, which are given in full.
%
% \ill{}: 0. \uncertain{}: 4 (all in the pencil marks, pp. 16 and 20).

\documentclass[11pt,a4paper]{article}
% Le préambule commun à toutes les transcriptions du fonds Grothendieck.
%
% Il tient en une idée : **l'incertitude se compose, elle ne se résout pas.**
% Une transcription qui lisse un mot douteux a détruit la seule chose pour
% laquelle on la faisait — un lecteur doit pouvoir savoir, à chaque endroit,
% ce qui était sur le feuillet et ce qui a été ajouté. Les six macros qui
% suivent sont donc l'appareil critique, et non de la décoration.
%
% Les mêmes commandes sont reconnues par `scripts/render.mjs`, qui produit la
% vue de lecture du site. Une macro ajoutée ici doit l'être aussi là-bas,
% faute de quoi le rendu échouera bruyamment — ce qui est voulu : un rendu
% silencieusement faux est pire qu'un rendu absent.

\ProvidesPackage{grothendieck}[2026/08/08 Transcriptions du fonds Grothendieck]

\usepackage{amsmath,amssymb,amsthm}
\usepackage{mathrsfs}
\usepackage{stmaryrd}
% Les diagrammes commutatifs sont partout dans ce fonds : c'est la langue
% même de la géométrie algébrique de Grothendieck, et une transcription qui
% les rendrait en prose ne transcrirait rien.
\usepackage{tikz-cd}
% Français par défaut — c'est la langue des feuillets — mais l'anglais est
% chargé aussi : la traduction et le résumé sont des documents anglais, et la
% typographie française (espaces avant les deux-points) y serait une faute.
% Ces documents basculent par \selectlanguage{english} en tête de corps.
\usepackage[english,french]{babel}
\usepackage{fontspec}
\usepackage{geometry}
\geometry{a4paper,margin=25mm,marginparwidth=28mm}
\usepackage{marginnote}
\usepackage{xcolor}
% ulem, not soul. `soul`'s \st and \ul are built on a character-by-character
% scan that XeLaTeX's font machinery breaks: the document fails with an
% undefined control sequence rather than degrading. ulem does the same two jobs
% and is Unicode-engine safe. `normalem` stops it hijacking \emph, which the
% transcriptions use for Grothendieck's own emphasis.
\usepackage[normalem]{ulem}
\usepackage{hyperref}
\hypersetup{colorlinks=true,linkcolor=black,urlcolor=black,pdfborder={0 0 0}}

% --- Métadonnées du lot -----------------------------------------------------
% Reprises telles quelles par le script de rendu, qui les affiche en tête de
% la vue de lecture. Elles fixent ce que le fichier prétend couvrir : sans
% elles, une transcription flotte sans point d'attache dans un fonds de seize
% mille feuillets.

\newcommand{\folder}[1]{\gdef\grothFolder{#1}}
\newcommand{\batch}[1]{\gdef\grothBatch{#1}}
\newcommand{\leaves}[2]{\gdef\grothFirst{#1}\gdef\grothLast{#2}}
\newcommand{\foldertitle}[1]{\gdef\grothFoldertitle{#1}}
\gdef\grothFolder{?}\gdef\grothBatch{?}\gdef\grothFirst{?}\gdef\grothLast{?}\gdef\grothFoldertitle{}

% --- Le résumé --------------------------------------------------------------
%
% Il ouvre la lecture modernisée, sous le titre « Résumé », et il est composé
% en petit. Ce n'est pas un ornement : c'est le seul passage du document qui
% ne soit pas des mathématiques, et un lecteur doit pouvoir le reconnaître
% avant de l'avoir lu — puis le sauter s'il connaît déjà le sujet, ou s'y
% arrêter s'il ne le connaît pas. Le filet qui le suit marque la reprise du
% corps.

\newenvironment{resume}
  {\small\setlength{\parskip}{0.45em}}
  {\par\medskip\hrule\medskip}

% --- Mots-clés ---------------------------------------------------------------
%
% En anglais, en clôture du résumé de la lecture modernisée : le vocabulaire
% moderne sous lequel chercher ce que les feuillets construisent. Le manifeste
% du site (`npm run manifest`) les extrait pour étiqueter la cote entière —
% cette ligne est la source unique des tags, il n'y a pas de fichier à côté.
\newcommand{\keywords}[1]{\par\smallskip\noindent
  {\footnotesize\textsc{Keywords} --- \textit{#1}}\par}

% --- L'appareil critique ----------------------------------------------------

% Le numéro de feuillet, en marge. C'est l'ancre qui permet de régler un
% désaccord sur une formule en regardant *un* feuillet plutôt qu'une chemise
% de six cents. Le site s'en sert aussi pour tourner les pages du fac-similé
% quand on fait défiler la transcription.
\newcommand{\leaf}[1]{%
  \marginnote{\footnotesize\color{gray!70}\textsf{#1}}%
  \label{leaf:#1}\ignorespaces}

% Illisible : on ne devine pas. Un mot inventé qui se lit comme les autres est
% le pire résultat possible.
\newcommand{\ill}{\textcolor{red!60!black}{[\,\ldots\,]}}

% Lecture douteuse : le mot est donné, et signalé comme incertain.
\newcommand{\uncertain}[1]{\uline{#1}}

% Ajout de l'éditeur — une lettre manquante, un mot sous-entendu.
\newcommand{\add}[1]{\textcolor{blue!50!black}{[#1]}}

% Biffé par Grothendieck lui-même. Ce qu'il a rayé fait partie du document :
% une rature dit souvent où la pensée a changé de direction.
\newcommand{\struck}[1]{\sout{#1}}

% Note du transcripteur, sur l'état matériel du feuillet.
\newcommand{\note}[1]{\par\smallskip{\small\color{gray!60!black}\textit{[#1]}}\par\smallskip}

% Note marginale de Grothendieck, distincte du corps du texte.
\newcommand{\marginal}[1]{\par\smallskip\noindent\hspace{1em}%
  {\small\color{gray!60!black}#1}\par\smallskip}

% --- Titre ------------------------------------------------------------------

\AtBeginDocument{%
  \begin{center}
    {\large\bfseries Fonds Alexandre Grothendieck --- cote n\textsuperscript{o}~\grothFolder}\\[2pt]
    {\itshape\grothFoldertitle}\\[4pt]
    {\small feuillets \grothFirst--\grothLast\ (lot \grothBatch)}\\[2pt]
    {\footnotesize\color{gray!60!black}Transcription établie d'après le fac-similé numérisé
     par l'Université de Montpellier.\\
     Les lectures incertaines sont soulignées, les illisibles notées [\,\ldots\,],
     les ajouts entre crochets.}
  \end{center}
  \vspace{1em}}

\endinput


\folder{63}
\batch{1}
\pages{1}{20}
\dating{[à partir de 1964-vers 1970]}
\watermark{Édition de démonstration}
\foldertitle{Formulaire courbes elliptiques : notes et copies de notes manuscrites (s.d.), tapuscrit (s.d.), copies d'articles (s.d.)}

\begin{document}

\page{2}
\note{Les pages 2 à 13 sont la copie d'un manuscrit qui n'est pas de la main de
Grothendieck : une écriture ronde et posée, en anglais, sans rien de sa main
rapide (comparer le dossier 62), et le texte dit lui-même « See Serre's
"Courbes Elliptiques - Formulaire" » (page 4) et « (Serre tells me) »
(page 8). C'est la lettre de J.~Tate à Cassels publiée, légèrement retouchée,
sous le titre « Algorithm for determining the type of a singular fiber in an
elliptic pencil », \emph{Modular Functions of One Variable IV}, Lecture Notes
in Math.~476 (1975), p.~33-52 ; la table ci-dessous reproduit exactement ses
sections 1 à 8. Texte d'un autre auteur, publié depuis : il n'est pas
recomposé ici, chaque page reçoit une note. Aucune annotation d'une autre main
n'y est visible.}
\note{Page 2 : table des matières (« Contents ») de Tate, renvoyant à sa propre
pagination : 1. Generalized Weierstrass form, p.~1 ; 2. Change of coordinates,
p.~2 ; 3. « Minimal » Weierstrass equations over a valuation ring, p.~3 ;
4. The canonical filtration on the group of $v$-adic points, p.~4 ;
5. Application : the relation between $L_v(1)$ and $\int_{E_v}|\omega_v|$ in
case $k$ is finite, p.~6 ; 6. The Néron minimum model, p.~7 ; 7. Algorithm for
analysing singular fibers (first five cases), p.~8 ; 8. Algorithm continued
(last five cases), p.~10.}

\page{3}
\note{Tate, p.~1 : §1 « Generalized Weierstrass form ». Équation (1.1)
$y^2+a_1xy+a_3y=x^3+a_2x^2+a_4x+a_6$ d'une courbe elliptique sur un corps $K$
munie d'un point rationnel $O$ ; les quantités $b_2,b_4,b_6,b_8$, $c_4,c_6$,
$\Delta$, $j=c_4^3/\Delta$ (1.2) et leurs relations (1.3)
$4b_8=b_2b_6-b_4^2$, $1728\Delta=c_4^3-c_6^2$ ; la différentielle de première
espèce $\omega$ (1.4)-(1.5) ; le passage (1.6)-(1.7) à la forme
$\eta^2=\xi^3-\tfrac{c_4}{48}\xi-\tfrac{c_6}{864}$ et la comparaison (1.8) avec
$\wp$, $g_2$, $g_3$.}

\page{4}
\note{Tate, p.~2 : §2 « Change of coordinates ». Les changements de
coordonnées (2.1) $x=u^2x'+r$, $y=u^3y'+su^2x'+t$, les formules (2.2) donnant
$a_i'$ et $b_i'$, les invariances (2.3) $u^4c_4'=c_4$, $u^6c_6'=c_6$,
$u^{12}\Delta'=\Delta$, $j'=j$ ; l'exemple (2.4) d'une courbe « générique » de
module $j$ ; renvoi au formulaire de Serre pour les automorphismes.}

\page{5}
\note{Tate, p.~3 : §3 « The "minimal" Weierstrass equations over a valuation
ring ». Définition 3.1 (équation minimale relativement à $v$), Théorème 3.2
(existence et unicité à un changement de coordonnées près avec $r,s,t\in R$,
$u$ inversible), Corollaire 3.3 (unicité de $\omega$ à une unité près) et
quatre remarques : critère $v(\Delta)<12$, discriminant minimal
$\mathfrak{D}_E=\sum v(\Delta_v)\,v$, cas d'un anneau principal, borne sur
$v(\Delta)$ lorsque $j\in R$.}

\page{6}
\note{Tate, p.~4 : §4 « The canonical filtration on the group of $v$-adic
points ». Réduction $\tilde E$, sa partie lisse $\tilde E_0$ (courbe
elliptique, groupe multiplicatif éventuellement tordu, ou groupe additif),
$E_0(K)=\rho^{-1}(\tilde E_0(k))$, Théorème 4.1 ($E_0(K)$ d'indice fini,
$\rho$ homomorphisme), les $E_m(K)$, Théorème 4.2 (paramètre uniformisant
$z=-x/y$, développements (4.3) à coefficients dans $R$, loi de groupe formel
$\Phi$).}

\page{7}
\note{Tate, p.~5 : démonstration du Théorème 4.2 par (4.4)-(4.6) (équation en
$w=-1/y$, $z$ ; somme des trois racines ; « (sauf erreur) » en marge, de la
main de Tate), puis la « canonical filtration » (4.7)
$E(K)\supset E_0(K)\supset E_1(K)\supset E_2(K)\supset\cdots$ et ses
quotients.}

\page{8}
\note{Tate, p.~6 : §5 « Application ». Corollaire 5.1 ($\int_{E_1(K)}|\omega|
=1/q$ pour $\omega$ issue d'une équation minimale), le facteur local (5.2)
$L_v(s)$ dans les trois cas de réduction, et Théorème 5.2 (5.3)
$\int_{E(K)}|\omega|=(E(K):E_0(K))/L_v(1)$ : les « fudge factors » de Birch et
Swinnerton-Dyer sont les indices $(E(K):E_0(K))$.}

\page{9}
\note{Tate, p.~7 : §6 « The Néron minimum model ». Le modèle régulier minimal
$X$ sur $R$, et le tableau des dix types de fibres (symboles de Kodaira
$\mathrm{I}_0$, $\mathrm{I}_\nu$, $\mathrm{II}$, $\mathrm{III}$, $\mathrm{IV}$,
$\mathrm{I}_0^*$, $\mathrm{I}_\nu^*$, $\mathrm{IV}^*$, $\mathrm{III}^*$,
$\mathrm{II}^*$ : dessins, nombre de composantes, groupe
$E(K)/E_0(K)$, $\tilde E_0(k)$, et, pour la caractéristique résiduelle
$\neq 2,3$, $v(\Delta)$, exposant du conducteur, comportement de $j$), puis
commentaires sur $\tilde X_0$ et la minimalité.}

\page{10}
\note{Tate, p.~8 : §7 « Algorithm for analysing singular fibers (first five
cases) ». Conducteur $\mathfrak{p}^f$ avec $f=v(\Delta)+1-n$, ordre $c$ de
$E(K)/E_0(K)$, groupe $\tilde E_0(k)$ ; étapes (1) (type $\mathrm{I}_0$) et
(2) (type $\mathrm{I}_\nu$, sous-cas (2a), (2b) selon le corps de
décomposition de $T^2+a_1T-a_2$).}

\page{11}
\note{Tate, p.~9 : étapes (3) (type $\mathrm{II}$), (4) (type $\mathrm{III}$,
équation (7.1) et conique (*)), (5) (type $\mathrm{IV}$), chacune avec sa
justification entre crochets.}

\page{12}
\note{Tate, p.~10 : §8 « Algorithm continued, the last five cases ». Le
polynôme (8.1) $P(T)=T^3+a_{2,1}T^2+a_{4,2}T+a_{6,3}$, l'équation (8.2) ;
étapes (6) (type $\mathrm{I}_0^*$) et (7) (type $\mathrm{I}_\nu^*$, avec la
« subprocedure » (7.1)-(7.2)) ; « (continued) » en bas de page.}

\page{13}
\note{Tate, p.~11 : fin de la sous-procédure de l'étape (7) (7.3), puis
étapes (8) (type $\mathrm{IV}^*$, (8.1)), (9) (type $\mathrm{III}^*$, (9.1))
et (10) (type $\mathrm{II}^*$ ; si $\mathfrak{p}^6\mid a_6$, l'équation
n'était pas minimale). Fin de la copie du manuscrit de Tate dans ce lot.}

\page{15}
\note{Les pages 15 à 20 sont les six premiers feuillets (pages dactylographiées
1 à 6) du tapuscrit « Courbes elliptiques : formulaire, d'après J.~Tate, mis
au goût du jour par P.~Deligne », publié sous ce titre dans \emph{Modular
Functions of One Variable IV}, Lecture Notes in Math.~476 (1975), p.~53-73.
Texte de Deligne et publié : il n'est pas recomposé ; chaque page reçoit une
note donnant en entier les marques manuscrites, toutes au crayon.}
\note{Page 1 du tapuscrit : titre, §1 « Formules générales » : courbes de genre
$1$ sur $S$, fibres de type elliptique, multiplicatif ou additif, faisceau
$\omega=e^*\Omega^1_{E/S}$. Au crayon : « genre arithmétique » est souligné
d'un trait ondulé, et un « ? » est porté dans la marge gauche en face.}

\page{16}
\note{Page 2 du tapuscrit : $R^1p_*\mathcal{O}(ne)=0$ pour $n>0$, les
$p_*\mathcal{O}(ne)$ localement libres de rang $n$, la filtration de gradué
(1.1), le faisceau dualisant $\Omega^{\mathrm{reg}}_{E/S}$. Au crayon, dans la
marge droite : \uncertain{rf} en face de « a un seul groupe de cohomologie
non nulle » ; $\omega_{E/S}$ en face de « Là où $p$ est lisse, on a » ;
\uncertain{rf} de nouveau en face de « de formation compatible à tout
changement de base » ; un petit trait oblique sépare « donc » de
« localement ».}

\page{17}
\note{Page 3 du tapuscrit : $p_*\Omega^{\mathrm{reg}}_{E/S}\simeq
e^*\Omega^1_{E/S}=\omega$ et $p^*\omega\simeq\Omega^{\mathrm{reg}}_{E/S}$,
différentielles invariantes, bases $(y,x,1)$ de $p_*\mathcal{O}(3e)$, les
changements (1.2), l'équation (1.3). Au crayon : un « ? » dans la marge droite
en face de $p^*\omega\simeq\Omega^{\mathrm{reg}}_{E/S}$ ; le « 1 » de
« $(x, y, 1)$ » est entouré d'une parenthèse.}

\page{18}
\note{Page 4 du tapuscrit : existence et unicité de l'équation, la forme
$\pi=dx/F_y$, la réciproque (cubique à point d'inflexion en $(0,1,0)$), les
formules (1.4) pour $b_2,b_4,b_6,b_8,d_6,d_8$. Au crayon : $2^2$ écrit
au-dessus du coefficient $4$ de $4a_2a_6$ ; $2\cdot 3^2$ au-dessus du
$18$ de $18b_6$ ; une accolade suivie de « ? » dans la marge droite en face de
« et présentant un point d'inflexion en $(0,1,0)$ » ; une parenthèse ouvrante
devant « $12\,d_6$ » et devant « $12\,d_8$ ».}

\page{19}
\note{Page 5 du tapuscrit : (1.5) $c_4$, $c_6$, $\Delta$ ; (1.6)-(1.7) les
formules de changement de coordonnées pour les $a_i'$, $b_i'$, $d_i'$. Au
crayon, les coefficients numériques de (1.5) sont décomposés en facteurs
premiers au-dessus de chacun : $2^3 3$ sur $24$, $2^2 3^2$ sur $36$, $2^3 3^3$
sur $216$ (entouré), $2^3$ sur $8$, $3^3$ sur $27$, $3^2$ sur $9$, puis dans
la seconde expression de $\Delta$, $3^3$ sur $27$ et $2^2 3^2$ sur $36$ ;
l'exposant de $b_2^2$ dans $b_8(36b_4-b_2^2)$ est entouré, avec un « ? »
au-dessus à droite.}

\page{20}
\note{Page 6 du tapuscrit : (1.8) $u^4c_4'=c_4$, $u^6c_6'=c_6$,
$u^{12}\Delta'=\Delta$ ; §2 « Cas où 2 et 3 sont inversibles » : (2.1)-(2.4),
retour à $Y^2=4X^3-g_2X-g_3$, et la Proposition 2.5 (schéma de modules
$\mathrm{Spec}\,\mathbf{Z}[2^{-1},3^{-1}][g_2,g_3]$ \ldots). Dans la marge,
au crayon, sous la phrase qui suit (1.8), une note dont l'écriture rapide
paraît être la sienne (attribution non certaine) :}

\marginal{\emph{\uncertain{Sur un corps}}
\begin{itemize}
  \item cas elliptique : $\Delta\neq 0$
  \item cas multiplicatif : $\Delta=0$, $c_4\neq 0$ (\uncertain{équiv.} à $c_6\neq 0$)
  \item cas additif : $\Delta=0$, $c_4=0$ i.e.\ $\Delta=0$, $c_6=0$
\end{itemize}}
\note{les trois cas sont réunis par une accolade ; « Sur un corps » est
souligné.}

\end{document}
